% Generated draft: requires mathematical and source review.
\documentclass[11pt]{article}
\usepackage[margin=25mm]{geometry}
\usepackage{fontspec}
\setmainfont{Latin Modern Roman}
\usepackage{amsmath,amssymb,mathtools,unicode-math}
\setmathfont{Latin Modern Math}
\usepackage{enumitem,needspace,xurl,hyperref}
\setlength{\parindent}{0pt}
\setlength{\parskip}{5pt}
\setlength{\emergencystretch}{3em}
\newcommand{\N}{\mathbb{N}}
\newcommand{\Z}{\mathbb{Z}}
\newcommand{\Q}{\mathbb{Q}}
\newcommand{\R}{\mathbb{R}}
\newcommand{\C}{\mathbb{C}}
\newcommand{\F}{\mathbb{F}}
\newcommand{\A}{\mathbb{A}}
\newcommand{\PP}{\mathbb{P}}
\newcommand{\E}{\mathbb{E}}
\newcommand{\eps}{\varepsilon}
\newcommand{\dd}{\,\mathrm{d}}
\begin{document}
\section*{UnsolvedMath ID: 701135}
\textbf{Source problem code:} catalogue entry
\subsection{Definitions and mathematical statement}
For positive integers \(n,k\), define \(\mathcal E(n,k)\) to mean that there exist a finite set \(\Omega\) of cardinality \(n\), permutation groups \(M<G\leq\operatorname{Sym}(\Omega)\), and a partition \(\mathcal P\) with \(k\) nonempty parts of
\[
\Omega^{(2)}=\{(\alpha,\beta)\in\Omega\times\Omega:\alpha\ne\beta\}
\]
satisfying the following conditions. The group \(M\) is transitive on \(\Omega\). Under the diagonal action \(g\cdot(\alpha,\beta)=(g\alpha,g\beta)\), the group \(G\) permutes the parts of \(\mathcal P\) transitively, and every element of \(M\) fixes every part setwise. Finally, the image of the induced homomorphism
\[
\rho:G\longrightarrow\operatorname{Sym}(\mathcal P),\qquad \rho(g)(P)=gP,
\]
is a noncyclic group. Equivalently, the parts describe a decomposition of the complete directed graph on \(\Omega\), without loops, into edge-disjoint isomorphic directed graphs, with the stated group actions.

For a prime \(r\) and a positive integer \(n\), define the \(r\)-part of \(n\) by
\[
n_r=r^{v_r(n)},\qquad v_r(n)=\max\{a\in\mathbb Z_{\geq0}:r^a\mid n\}.
\]
Write \(C(n,k)\) for the condition
\[
 n_r\equiv1\pmod{k}\qquad\text{for every prime }r.
\]
In particular, \(n_r=1\) when \(r\nmid n\). The supplied source reports that, under the same hypotheses on \(M,G,\Omega,\mathcal P\), cyclicity of \(\rho(G)\) implies \(C(n,k)\), and attributes this result to the cited 2003 paper.

The problem asks for examples with noncyclic induced group for values \(n,k\) failing \(C(n,k)\). The existential reading is
\[
\exists n,k\in\mathbb Z_{>0}\quad
\bigl(\mathcal E(n,k)\ \text{and}\ \exists\text{ a prime }r:\ n_r\not\equiv1\pmod{k}\bigr).
\]
\textit{Unresolved quantifier qualification.} The source asks for examples ``for any \(n,k\)'' not satisfying the congruence condition. It does not explicitly distinguish the existential reading above from a request for examples for every prescribed pair of parameters violating the condition. Under the latter reading, the intended range of admissible pairs also requires clarification; the source supplies no such range. Both readings are retained for review.
\subsection{Short English statement}
Seek a transitive permutation group \(M\) properly contained in \(G\), together with a partition of all ordered pairs of distinct points that \(M\) fixes part by part and \(G\) permutes transitively. Can the induced group on the parts be noncyclic while some prime-power part \(n_r\) of the number of points fails \(n_r\equiv1\pmod{k}\), where \(k\) is the number of parts? The source's phrase ``for any \(n,k\)'' leaves unresolved whether one example or a family covering prescribed parameters is requested.
\subsection{Sources}
\begin{itemize}
\item Local UnsolvedMath registry, permanent ID \(701135\), exact source problem code \texttt{catalogue entry}, title \emph{Kourovka Notebook Problem 15.47}. The supplied record labels the problem open; this is a supplied status, not an independently verified mathematical or literature assessment.
\item \emph{The Kourovka Notebook (Unsolved Problems in Group Theory)}, original problem 15.47, 15th issue (2002); proposed by C. H. Li and C. E. Praeger. The supplied background also identifies the 21st issue edition (2026), \texttt{arXiv:1401.0300v48}. The collection description names E. I. Khukhro and V. D. Mazurov as editors. These bibliographic details are supplied attribution and have not been independently checked.
\item The problem statement cites \emph{Transactions of the American Mathematical Society}, volume 355, number 2 (2003), pages 637--653, for the congruence condition when the induced group on the parts is cyclic. It describes the result as work of the proposers. No article title, separate author list, or link is supplied.
\end{itemize}
\end{document}
